% ============================================================
% COMANDOS PERSONALIZADOS - prjt-tcc-unesp-rc
% ============================================================

% ----- CONTROLE DE ATIVAÇÃO DOS PRÉ-TEXTUAIS -----
\newcommand{\exibirResumo}{%
    % ============================================================
% RESUMO
% ============================================================

% ----- PALAVRAS-CHAVE (EDITAR AQUI) -----
\newcommand{\palavrachaveum}{Estudo dirigido}
\newcommand{\palavrachavedois}{Ensino de Geografia}
\newcommand{\palavrachavetres}{Currículo Paulista}
\newcommand{\palavrachavequatro}{Formação de Professores}
\newcommand{\palavrachavecinco}{} % Deixe vazio se não houver

\begin{resumo}

% --- TEXTO DO RESUMO (12pt, ESPAÇAMENTO SIMPLES, SEM RECUO, JUSTIFICADO) ---
\setlength{\parindent}{0pt}
\setlength{\parskip}{0pt}
\singlespacing
\normalsize

Escreva aqui o texto do seu resumo. Ele deve ser um parágrafo único, conter entre 150 e 500 palavras e apresentar de forma clara a introdução, o objetivo, a metodologia, os resultados esperados e as conclusões do seu projeto.

\vspace{0.5cm}

% --- FORMATAÇÃO AUTOMÁTICA DAS PALAVRAS-CHAVE ---
{\normalsize\bfseries Palavras-chave:} \normalsize
\palavrachaveum. \palavrachavedois. \palavrachavetres.%
\ifx\palavrachavequatro\empty\else
    \ \palavrachavequatro.%
\fi
\ifx\palavrachavecinco\empty\else
    \ \palavrachavecinco.%
\fi

\end{resumo}

% ============================================================
% FIM DO ARQUIVO resumo.tex
% ============================================================%
}

\newcommand{\exibirAbstract}{%
    % ============================================================
% ABSTRACT
% ============================================================

% ----- TÍTULO E SUBTÍTULO EM INGLÊS (EDITAR AQUI) -----
\newcommand{\tituloingles}{Title in English}
\newcommand{\subtituloingles}{Subtitle in English (optional)} % Deixe vazio se não houver

% ----- KEYWORDS (EDITAR AQUI) -----
\newcommand{\keywordum}{Keyword 1}
\newcommand{\keyworddois}{Keyword 2}
\newcommand{\keywordtres}{Keyword 3}
\newcommand{\keywordquatro}{} % Deixe vazio se não houver
\newcommand{\keywordcinco}{} % Deixe vazio se não houver

\begin{abstract}

% --- TEXTO DO ABSTRACT (12pt, ESPAÇAMENTO SIMPLES, SEM RECUO, JUSTIFICADO) ---
\setlength{\parindent}{0pt}
\setlength{\parskip}{0pt}
\singlespacing
\normalsize

Write your abstract here. It should be a single paragraph between 150 and 500 words, clearly presenting the introduction, objective, methodology, expected results, and conclusions of your project.

\vspace{0.5cm}

% --- FORMATAÇÃO AUTOMÁTICA DAS KEYWORDS ---
{\normalsize\bfseries Keywords:} \normalsize
\keywordum. \keyworddois. \keywordtres.%
\ifx\keywordquatro\empty\else
    \ \keywordquatro.%
\fi
\ifx\keywordcinco\empty\else
    \ \keywordcinco.%
\fi

\vspace{0.5cm}

% --- TÍTULO EM INGLÊS (Mesma formatação das Keywords) ---
{\normalsize\bfseries Title:} \normalsize
\tituloingles%
\ifx\subtituloingles\empty\else
    : \subtituloingles%
\fi

\end{abstract}

% ============================================================
% FIM DO ARQUIVO abstract.tex
% ============================================================%
    \clearpage
}

% ----- COMANDO PARA EXIBIR SUMÁRIO (SEM NUMERAÇÃO) -----
\newcommand{\exibirSumario}{%
    \vspace*{0pt}
    \begin{center}
        {\fontsize{14pt}{16pt}\selectfont\bfseries SUMÁRIO}
    \end{center}
    \vspace{-3cm}
    \renewcommand{\contentsname}{}
    \thispagestyle{empty}
    \tableofcontents
    \thispagestyle{empty}
    \clearpage
}

% ----- COMANDOS PARA PÓS-TEXTUAIS -----
\newcommand{\exibirApêndices}{%
    \ifnum\ativarApêndices=1
        % ============================================================
% APÊNDICES - trab-disc-unesp
% ============================================================

\chapter{APÊNDICE A - Título do Apêndice}
\label{apendice:a}

Texto do apêndice...

\chapter{APÊNDICE B - Título do Apêndice}
\label{apendice:b}

Texto do apêndice...

% ============================================================
% FIM DO ARQUIVO 02-apendices.tex
% ============================================================%
    \fi
}

\newcommand{\exibirAnexos}{%
    \ifnum\ativarAnexos=1
        % ============================================================
% ANEXOS - trab-disc-unesp
% ============================================================

\chapter{ANEXO A - Título do Anexo}
\label{anexo:a}

Texto do anexo...

\chapter{ANEXO B - Título do Anexo}
\label{anexo:b}

Texto do anexo...

% ============================================================
% FIM DO ARQUIVO 01-anexos.tex
% ============================================================%
    \fi
}

% ----- AMBIENTES PRÉ-TEXTUAIS -----
\newenvironment{resumo}{%
    \thispagestyle{empty}
    \vspace*{0pt}
    \begin{center}
        {\fontsize{14pt}{16pt}\selectfont\bfseries RESUMO}
    \end{center}
    \vspace{0.5cm}
}{%
    \par
}

\newenvironment{abstract}{%
    \thispagestyle{empty}
    \vspace*{0pt}
    \begin{center}
        {\fontsize{14pt}{16pt}\selectfont\bfseries ABSTRACT}
    \end{center}
    \vspace{0.5cm}
}{%
    \par
}

% ============================================================
% FORMATAÇÃO DE CAPÍTULOS E SEÇÕES (MÉTODO INFALÍVEL)
% ============================================================

% 1. Remove qualquer formatação anterior de capítulos
\titleformat{\chapter}[block]{}{}{}{}
\titlespacing*{\chapter}{0pt}{0pt}{0pt}

% 2. Cria o comando para o título do capítulo (centralizado, 14pt, negrito, maiúsculo)
\newcommand{\meuTituloCapitulo}[1]{%
    \vspace{1.5cm} % Espaço antes do capítulo
    \begin{center}
        {\fontsize{14pt}{16pt}\selectfont\bfseries \MakeUppercase{#1}}
    \end{center}
    \vspace{0.5cm} % Espaço depois do capítulo
}

% 3. Redefine o comando \chapter usando a classe book nativa
\makeatletter
\renewcommand{\chapter}[1]{%
    \par\refstepcounter{chapter}%
    \meuTituloCapitulo{#1}%
    \addcontentsline{toc}{chapter}{\protect\numberline{\thechapter}#1}%
}
\makeatother

% 4. Configuração das seções (esquerda, negrito)
\titleformat{\section}
    {\normalfont\normalsize\bfseries\raggedright}
    {\thesection}
    {0.5em}
    {}

\titleformat{\subsection}
    {\normalfont\normalsize\bfseries\raggedright}
    {\thesubsection}
    {0.5em}
    {}

\titleformat{\subsubsection}
    {\normalfont\normalsize\itshape\raggedright}
    {\thesubsubsection}
    {0.5em}
    {}

% ============================================================
% CONFIGURAÇÃO DO CABEÇALHO (NUMERAÇÃO SUPERIOR DIREITA)
% ============================================================
\pagestyle{fancy}
\fancyhf{}
\fancyhead[R]{\thepage}
\renewcommand{\headrulewidth}{0pt}
\renewcommand{\chaptermark}[1]{}
\renewcommand{\sectionmark}[1]{}

\fancypagestyle{plain}{%
    \fancyhf{}%
    \fancyhead[R]{\thepage}%
    \renewcommand{\headrulewidth}{0pt}%
}

% ============================================================
% COMANDOS ÚTEIS PARA FIGURAS E TABELAS
% ============================================================
\newcommand{\figura}[4]{%
    \begin{figure}[H]
        \centering
        \includegraphics[width=#1]{#2}
        \caption{#3}
        \label{fig:#4}
    \end{figure}
}

\newcommand{\tabela}[4]{%
    \begin{table}[H]
        \centering
        \caption{#1}
        \label{tab:#2}
        #3
        \fonte{#4}
    \end{table}
}

\newcommand{\fonte}[1]{%
    \vspace{0.2cm}
    \par\small\textbf{Fonte:} #1
}

\newcommand{\nota}[1]{%
    \vspace{0.2cm}
    \par\small\textbf{Nota:} #1
}

\newcommand{\destaque}[1]{\textbf{#1}}
\newcommand{\citacao}[1]{\textit{#1}}

% ============================================================
% FIM DO ARQUIVO comandos.tex
% ============================================================